% ch-a4-12.tex —— 第 4 篇（Emerton--Gee--Hellmann）续片，底本 PDF 312--323 = 印刷页 299--310
\ysrc{299}
与主序列情形一样，我们取系数域 $L$ 在 $\Q_p$ 上非分歧（驯顺型总可以定义在非分歧扩张上）；在该假设下我们看到，$\XX^{\mathrm{crys},(1,0),\tau}$ 仍是正则的，从而层 $\mathfrak A\bigl((c\text{-}\mathrm{Ind}_K^G\Lambda_i)^\wedge\bigr)$ 仍是可逆的。此外，由于 $\Lambda_1^\vee=\Lambda_2$，我们发现 $\mathfrak A\bigl((c\text{-}\mathrm{Ind}_K^G\Lambda_1)^\wedge\bigr)$ 与 $\mathfrak A\bigl((c\text{-}\mathrm{Ind}_K^G\Lambda_2)^\wedge\bigr)$ 在 Grothendieck--Serre 对偶下必然互换。

现在拉回给出同构
\[
\mathrm{Pic}\Bigl(\wh{\PP}^{1}\times\bigl[(\Spf(\OO[x,y]/(xy-p))^{\wedge}/\wh{\GG}_m)\bigr]\Bigr)
\isom\mathrm{Pic}\bigl(\XX^{\mathrm{crys},(1,0),\tau}\bigr),
\]
于是 $\XX^{\mathrm{crys},(1,0),\tau}$ 上任一可逆层都必然是 $\mathcal L(n)$ 的形式，其中 $\mathcal L$ 由 $\wh{\PP}^{1}$ 拉回，$n$ 是表明 $\wh{\GG}_m$-作用之权的整数。此外，$\XX^{\mathrm{crys},(1,0),\tau}$ 上的对偶层由 $\wh{\PP}^{1}\times\bigl[(\Spf(\OO[x,y]/(xy-p))^{\wedge}/\wh{\GG}_m)\bigr]$ 上的对偶层 $\OO(-2)_{\wh{\PP}^{1}}$ 拉回得到。

短正合列
\[
0\to(c\text{-}\mathrm{Ind}_K^G\Lambda_2)^\wedge\to(c\text{-}\mathrm{Ind}_K^G\Lambda_1)^\wedge\to c\text{-}\mathrm{Ind}_K^G\sigma\to0
\]
与
\[
0\to(c\text{-}\mathrm{Ind}_K^G\Lambda_1)^\wedge\xrightarrow{\ p\ }(c\text{-}\mathrm{Ind}_K^G\Lambda_2)^\wedge\to c\text{-}\mathrm{Ind}_K^G\sigma^{\mathrm{co}}\to0
\]
在施加 $\mathfrak A$ 之后诱导出相应的层短正合列；而由 $\mathfrak A$ 在 $c\text{-}\mathrm{Ind}_K^G\sigma$ 与 $c\text{-}\mathrm{Ind}_K^G\sigma^{\mathrm{co}}$ 上的已知取值，我们得到
\[
\mathfrak A\bigl((c\text{-}\mathrm{Ind}_K^G\Lambda_1)^\wedge\bigr)\cong\mathcal L(-1),
\]
以及
\[
\mathfrak A\bigl((c\text{-}\mathrm{Ind}_K^G\Lambda_2)^\wedge\bigr)\cong\mathcal L(1),
\]
并且态射 (7.4.7) 由乘以 $y$ 诱导。这一分析并不确定可逆层 $\mathcal L$（因为从 $\PP^{1}$ 拉回的任何 $\mathcal L$，无论限制到 $\XX(\sigma)$ 还是 $\XX(\sigma^{\mathrm{co}})$ 上，都必然平凡）。然而 $\mathcal L(\pm1)$ 的相互对偶性随后表明，$\mathcal L$ 必是 $\OO(-1)_{\wh{\PP}^{1}}$ 的拉回。

\begin{rem}{7.4.8}\label{thm:a4-7-4-8}
尖点情形与主序列情形的一个差别在于：在尖点情形，我们有 $\Gamma\bigl(\XX^{\mathrm{crys},(1,0),\tau},\OO_{\XX^{\mathrm{crys},(1,0),\tau}}\bigr)=\OO$（正如 $\Gamma\bigl(\XX(\sigma|\sigma^{\mathrm{co}}),\OO_{\XX(\sigma|\sigma^{\mathrm{co}})}\bigr)=k$）。经由对函子 $\mathfrak A$ 的考察，这反映了如下事实：紧诱导表示 $\bigl(c\text{-}\mathrm{Ind}_K^G\sigma^{\mathrm{crys},\circ}(\tau)\bigr)^\wedge$ 的自同态环（对晶格 $\sigma^{\mathrm{crys},\circ}(\tau)$ 的任一选取）在此情形下是平凡的。
\end{rem}

\subsubsection*{7.4.9.\quad 与解析情形的相容性示例}
下面暂时允许 $\tau$ 表示驯顺主序列型或驯顺尖点型。此时，正如半稳定情形中 (5.2.14) 所指出的，存在态射
\[
\pi:(\XX^{\mathrm{crys},(1,0),\tau})^{\mathrm{rig}}_\eta\to(\mathrm{Fil}_{(1,0)}\mathrm{WD}_\tau)^{\mathrm{an}},
\]
其目标是二维 Weil--Deligne 表示的（刚性解析）模叠：这些表示的惯性型为 $\tau$、逆分圆行列式固定，并带有一维子空间给出的滤过。

现设 $\tau$ 为主序列型。则
\begin{align*}
(\XX^{\mathrm{crys},(1,0),\tau})^{\mathrm{rig}}_\eta
&\cong\Bigl[\Bigl(\{(t,x)\in(\GG_m)_t\times\PP^{1}_x\mid 1\ge|t|\ge|p|\}\\
&\qquad\setminus\bigl(\{(x,t)\mid 1>|t|,|x|\}\cup\{(x,t)\mid |t|>|p|,|x|>1\}\bigr)\Bigr)\\
&\qquad\Big/\bigl(\wh{\GG}_m\bigr)^{\mathrm{rig}}_\eta\Bigr],
\end{align*}
其中\ysrc{300}$(\wh{\GG}_m)^{\mathrm{rig}}_\eta=\{u\mid |u|=1\}\subset\GG_m$，作用为 $u\cdot(t,x)=(t,u^2x)$；而\footnote{这里我们把乘法群与射影直线（视作刚性解析空间）径直记作 $\GG_m$ 与 $\PP^{1}$。}
\[
(\mathrm{Fil}_{(1,0)}\mathrm{WD}_\tau)^{\mathrm{an}}\cong[(\GG_m)_t\times\PP^{1}_x/(\GG_m)_u],
\]
其作用同样由 $u\cdot(t,x)=(t,u^2x)$ 给出。映射 $\pi$ 此时即是显然的那个。\tn{底本主序列与尖点情形的这两个公式中，集合 $\{(x,t)\mid 1>|t|,|x|\}$ 的 $|x|$ 之后均疑脱漏比较条件；译文照录。}

按此描述，「忘掉滤过」映射 $\mathrm{pr}_{WD}$ 是态射
\[
[(\GG_m)_t\times\PP^{1}_x/(\GG_m)_u]\to[(\GG_m)_t/(\GG_m)_u]
\]
（右端的叠是就 $(\GG_m)_u$ 在 $(\GG_m)_t$ 上的平凡作用构成的）。出现在 (6.2.5) 中的扭 $((-\rho')_\sigma)$ 在此情形下就是如下扭：沿「忘掉 Weil--Deligne 表示」映射
\[
\mathrm{pr}_{Fil}:[(\GG_m)_t\times\PP^{1}_x/(\GG_m)_u]\to\PP^{1}_x/\mathrm{SL}_2
\]
拉回层 $\OO(-1)_{\PP^{1}}$ 所得的扭；该映射由到第二个因子的投影以及 $(\GG_m)_u$ 作为 $\mathrm{SL}_2$ 中对角环面的嵌入共同诱导。\tn{底本此处作「along the the ...」，衍一 the，疑为笔误。}

现在，光滑范畴 Langlands 对应把 $c\text{-}\mathrm{Ind}_K^G\sigma^{\mathrm{crys}}(\tau)$ 对应到 $[(\GG_m)_t/(\GG_m)_u]$ 上带有平凡 $(\GG_m)_u$-等变结构的平凡可逆层。于是 §6.2.11 中预期的相容性 (6.2.13)（同时计入猜想 6.2.4 (1)）预测
\[
\mathfrak A\bigl((c\text{-}\mathrm{Ind}_K^G\sigma^{\mathrm{crys},\circ}(\tau))^\wedge\bigr)^{\mathrm{rig}}_\eta
\cong\mathrm{pr}_{Fil}^{*}\bigl(\OO(-1)_{\PP^{1}}\bigr).
\]
$\PP^{1}_x$ 上的线丛 $\OO(-1)_{\PP^{1}}$ 在投影到 $\PP^{1}_x$ 之下拉回到 $(\GG_m)_t\times\PP^{1}_x$ 的开子集
\begin{align*}
&\{(t,x)\mid 1\ge|t|\ge|p|\}\\
&\qquad\setminus\bigl(\{(x,t)\mid 1>|t|,|x|\}\cup\{(x,t)\mid |t|>|p|,|x|>1\}\bigr)
\end{align*}
上时成为平凡线丛，但它带有非平凡的 $(\GG_m)_u$-作用：事实上 $(\GG_m)_u$ 以权 $-1$ 作用于其上。因此 (6.2.13) 预测同构
\[
\mathfrak A\bigl((c\text{-}\mathrm{Ind}_K^G\sigma^{\mathrm{crys},\circ}(\tau))^\wedge\bigr)^{\mathrm{rig}}_\eta
\cong\OO_{(\XX^{\mathrm{crys},(1,0),\tau})^{\mathrm{rig}}_\eta}(-1),
\]
这与 §7.4.5 的计算相符。

现设 $\tau$ 为尖点型。此时
\begin{align*}
(\XX^{\mathrm{crys},(1,0),\tau})^{\mathrm{rig}}_\eta
&\cong\Bigl[\Bigl(\{(t,x)\in\PP^{1}\times\GG_m\mid 1\ge|x|\ge|p|\}\\
&\qquad\setminus\bigl(\{(t,x)\mid 1>|t|,|x|\}\cup\{(t,x)\mid |t|>1,|x|>|p|\}\bigr)\Bigr)\\
&\qquad\Big/\bigl(\wh{\GG}_m\bigr)^{\mathrm{rig}}_\eta\Bigr],
\end{align*}
其中 $(\wh{\GG}_m)^{\mathrm{rig}}_\eta$ 经 $u\cdot(t,x)=(t,u^2x)$ 作用；而
\[
(\mathrm{Fil}_{(1,0)}\mathrm{WD}_\tau)^{\mathrm{an}}\cong[\PP^{1}_t/\mu_2],
\]
其中 $\mu_2$ 平凡地作用。映射 $\pi$ 此时即是显然的那个，由忘掉 $x$ 给出。

映射 $\mathrm{pr}_{WD}$ 是投影
\[
[\PP^{1}_t/\mu_2]\to[*/\mu_2],
\]
而我们又有映射
\[
\mathrm{pr}_{Fil}:[\PP^{1}_t/\mu_2]\to[\PP^{1}_t/\mathrm{SL}_2]
\]
\ysrc{301}该映射忘掉 Weil--Deligne 表示，并由 $\mu_2$ 到 $\mathrm{SL}_2$ 的中心嵌入诱导。于是预期的相容性 (6.2.13) 预测同构
\[
\mathfrak A\bigl((c\text{-}\mathrm{Ind}_K^G\sigma^{\mathrm{crys},\circ}(\tau))^\wedge\bigr)^{\mathrm{rig}}_\eta
\cong\bigl(\pi^{*}\mathrm{pr}_{Fil}^{*}\OO(-1)_{\PP^{1}}\bigr)(-1)
\]
（外层的扭 $-1$ 表明 $(\wh{\GG}_m)^{\mathrm{rig}}_\eta$ 的等变作用），这与 §7.4.6 的计算相符。（注意：扭 $1$ 与扭 $-1$ 给出 $\XX^{\mathrm{crys},(1,0),\tau}$ 上互不同构的可逆层，但这些线丛拉回到 $(\XX^{\mathrm{crys},(1,0),\tau})^{\mathrm{rig}}_\eta$ 之后变得同构，因为可逆层 $\OO_{(\XX^{\mathrm{crys},(1,0),\tau})^{\mathrm{rig}}_\eta}(2)$ 平凡——它确实以坐标 $x$ 为处处非零的整体截面。）

\subsection{半正交分解}
我们现在阐述 $\IndCoh(\XX)$ 的一个显式半正交分解（其中 $\XX$ 如 §7.2），它使我们能描述函子 $\mathfrak A$ 的本质像。我们还要讨论 $D^\times$（$D$ 是以 $\Q_p$ 为中心的非分裂四元数代数）的表示以及 $\mathrm{GL}_2(\Q_{p^2})$ 的表示所对应的预期函子的类似本质像。

首先注意，我们有态射 $f:\mu_2\times\XX\to\mathcal I_{\XX}$（$\XX$ 的惯性叠），它由 $\zeta\in\mu_2$ 在 $(\varphi,\Gamma)$-模上乘以 $\zeta$ 的作用给出（注意由 $\zeta^2=1$ 知该作用在 $\wedge^2D$ 上平凡）。这个态射还（由构造可见）关于 $\mathcal I_{\XX}$ 上的群胚结构是中心的（$\mathcal I_{\XX}$ 是 $\XX$ 上的群胚）。

态射 $f$ 诱导 $\mu_2$ 在 $\Coh(\XX)$ 上、从而在 $\IndCoh(\XX)$ 上的一个典范作用，如下所述：取形式概形 $U$ 与光滑满射 $U\to\XX$（换言之，$\XX$ 的一张光滑图册），并令 $R:=U\times_{\XX}U$，于是可以写 $\XX=[U/R]$。态射 $f$ 随后在 $U$ 上诱导群胚的态射 $U\times\mu_2\to R$，其中 $\mu_2$ 在 $U$ 上平凡地作用。$\XX$ 上的凝聚层可以看作 $U$ 上带有 $R$-等变结构的凝聚层（这编码了下降到 $\XX$ 的下降数据）；特别地，任何这样的层都带有 $\mu_2$-等变结构，而由 $f$ 的中心性，该结构与 $R$-等变结构交换。因此 $\XX$ 上的凝聚层确实带有典范的 $\mu_2$-等变结构。

令 $\IndCoh(\XX)^{\pm}$ 是 $\IndCoh(\XX)$ 中 $-1\in\mu_2$ 以 $\pm$ 作用的子范畴。我们有时称 $\IndCoh(\XX)^{+}$ 为偶子范畴、$\IndCoh(\XX)^{-}$ 为奇子范畴。

作为走向 $\IndCoh(\XX)$ 的半正交分解的第一步，容易看出我们有正交分解，分成 $\IndCoh(\XX)^{+}$ 与 $\IndCoh(\XX)^{-}$ 两块。由 $L_\infty$ 的构造可得 $L_\infty$ 是奇的，因此 $\mathfrak A$ 的本质像包含在 $\IndCoh(\XX)^{-}$ 中。

\subsubsection*{7.5.1.\quad 动机}
关于半正交分解的一般性讨论，我们请读者参阅附录 A.8。我们的半正交分解由如下可能性所驱动：即注 1.4.6 中所谓的「$p$-进 Fargues--Scholze 猜想」；同时也由 $G$-丛的叠 $\Bun_G$ 在 Fargues--Fontaine 曲线上的结构所驱动，此处 $G=\mathrm{GL}_2(F)$。该叠容许一个分层（stratification），其分层片与特定的等晶体（isocrystal）相对应；而且正如 $\ell$-进情形一样，可以合理地期望：无论考虑 $\Bun_G$ 上何种 $p$-进层的范畴，都会有相应的半正交分解，分成若干块，每一块都可用等晶体的自同构群的表示论来描述。我们\ysrc{302}要提醒读者，我们并不预期本节所考虑的 $\IndCoh(\XX)$ 上的半正交分解会与来自 $\Bun_G$ 的（预期中的）半正交分解精确相符；相反，我们预期二者由一个「上三角幂幺基变换矩阵」相联系，而这个矩阵可能相当复杂。\footnote{Xinwen Zhu 曾向作者们指出，这一情形在某种程度上类似于 [Yun09, §4] 所考虑的层论 Radon 变换。}

\begin{rem}{7.5.2}\label{thm:a4-7-5-2}
把 $p$-进局部 Langlands 表述为范畴等价时，很可能需要对 $\IndCoh(\XX)$ 加上幂零奇异支撑（nilpotent singular support）条件。本节通篇忽略这一点；注意我们大部分时候只涉及叠 $\XX$ 的一般部分，在那里预期该条件是自动成立的。
\end{rem}

\subsubsection*{7.5.3.\quad 来自 $\Bun_G$ 的考量}
正如已经指出的，我们想要展示的 $\IndCoh(\XX)$ 的半正交分解，其结构由一个平行的结构所驱动；后者存在于 $\Bun_G$ 上的可构造层范畴中（其中 $G=\mathrm{GL}_2(F)$ 是我们关心的情形），来自 $\Bun_G$ 上的一个自然分层，我们现在来回忆它。

具体地，依照 [FS24, §I.4] 与 [FS25, §2.6]，叠 $\Bun_{\mathrm{GL}_2}$ 容许如下分层，分成局部闭的分层片。$\Bun_{\mathrm{GL}_2}$ 的连通分支由 $\alpha\in\frac{1}{2}\Z$ 标号，且每个连通分支包含唯一一个开分层片。若 $\alpha\in\Z$，该开分层片同构于 $[*/\mathrm{GL}_2(F)]$；而若 $\alpha\notin\Z$，它同构于 $[*/D^\times]$，其中 $D$ 是以 $F$ 为中心的非分裂四元数代数。（下文我们把群概形的 $p$-进点与对应的层不正式地混为一谈。）其余的分层片由整数对 $(m,n)$ 标号，其中 $m>n$ 且 $m+n=2\alpha$，并等同于商叠 $[*/\mathbf G_{m,n}]$，其中
\[
\mathbf G_{m,n}=\begin{pmatrix}F^\times&\mathrm{BC}(\OO(m-n))\\0&F^\times\end{pmatrix}
\]
是幂幺群菱形（实际上是 Banach--Colmez 空间）$\mathrm{BC}(\OO(m-n))$ 与环面 $F^\times\times F^\times$ 的半直积。（叠 $\Bun_{\mathrm{PGL}_2}$ 可由类似的数据描述：把相差一个由 $\OO(1)$ 给出的扭的等晶体等同起来；例如其连通分支现在由 $\alpha\in\frac{1}{2}\Z/\Z$ 标号。）

在 $\ell$-进情形，$\Bun_{\mathrm{GL}_2}$ 上的层范畴有相应的（无穷）半正交分解，见 [FS24, Thm. I.5.1]；我们预期 $p$-进情形也同样如此。（事实上，我们下文展示的 $\IndCoh(\XX)$ 的半正交分解正是由这一预期所驱动，反过来也印证了它。）在 $\ell$-进情形还可以忽略 Banach--Colmez 空间 $\mathrm{BC}(\OO(m-n))$，但对 $p$-进系数而言这一点先验地并不清楚。如下文所解释的，我们预期它们影响我们的半正交分解当且仅当 $[F:\Q_p]\ge m-n$。（特别地，对 $\mathrm{GL}_2(\Q_p)$，我们发现它们只产生 $\IndCoh(\XX)^{+}$ 的半正交分解中的单个成分，而在 $\IndCoh(\XX)^{-}$ 的分解中完全不起作用——因此特别地也不出现在 $\mathfrak A$ 的本质像的描述中。）

正如已经指出的，在 $F=\Q_p$ 的情形，我们将在下文构造 $\IndCoh(\XX)$ 的一个半正交分解，并特别地把 $\mathfrak A$ 的像描述为如下范畴的左正交补：该范畴由从 $\XX$ 中的\ysrc{303}可约轨迹推出的层所生成。我们进一步给出可约轨迹中的层的半正交分解，具体说是剩余可约轨迹上的层。如上所述，我们预期后一分解最终应与 $\Q_p^\times$ 的表示论相关联。此外，它应当可以用一个「谱 Eisenstein 函子」来描述，其中涉及注 5.3.16 中 Borel $(\varphi,\Gamma)$-模的叠在 Banach 情形下的类似物。

\subsubsection*{7.5.4.\quad 限制到可约轨迹}
我们回到 $F=\Q_p$ 的设定。剩余可约轨迹 $\mathcal U\subset\XX$ 是连通分支 $\mathcal U(\sigma|\sigma^{\mathrm{co}})$ 的并；为叙述方便，一如既往地假定 $\sigma$ 非常一般（very generic）。如 (7.2.23) 所见，我们有 $\mathcal U(\sigma|\sigma^{\mathrm{co}})=[\Spf\,\OO[s,t^{\pm1},x,y]^{\wedge}/\wh{\GG}_m]$，其中上标楔形表示 $(p,s,xy)$-进完备化，且 $x$ 与 $y$ 的权分别为 $2$ 与 $-2$。

令 $A:=\OO[s,t^{\pm1}]$，并令 $B=A[x,y]$，把它看作分次 $A$-代数，其中 $x$ 的权为 $2$、$y$ 的权为 $-2$。令 $I:=(p,s)\subset A$ 与 $J:=(I,xy)\subset B$ 为所指的理想，而 $\mathcal B_J$ 记分次 $B$-模范畴的全子范畴，由那些每个元素都被 $J$ 的某个幂零化的模组成。于是有典范等价
\[
\IndCoh\bigl(\mathcal U(\sigma|\sigma^{\mathrm{co}})\bigr)\isom\IndCoh(\mathcal B_J)
\]
（右端的定义见 (A.6.8) 以及 §C.4.16 的讨论），它在源与靶上都保持 $\pm$-分解（源上关于 $\mu_2$-作用，靶上关于分次的奇偶性）。

由命题 C.4.18 立即可得，每个 $\pm$-成分 $\IndCoh(\mathcal B_J^{\pm})$ 都容许半正交分解，分成半正交的块 $\mathcal D_0^{\pm}$ 与 $\mathcal E_0^{\pm}=(\mathcal D_0^{\pm})^{\perp}$，其中
\[
\mathcal D_0^{+}=\langle B/J\rangle,
\]
\[
\mathcal E_0^{+}=\bigl\langle(A/I)[x](-2),(A/I)[x](-4),\cdots,(A/I)[y](2),(A/I)[y](4),\cdots\bigr\rangle,
\]
\[
\mathcal D_0^{-}=\langle B/J(1),B/J(-1)\rangle,
\]
以及
\[
\mathcal E_0^{-}=\bigl\langle(A/I)[x](-3),(A/I)[x](-5),\cdots,(A/I)[y](3),(A/I)[y](5),\cdots\bigr\rangle.
\]
由 §7.3.20 中我们对函子 $\mathfrak A$ 限制到 $\mathcal U(\sigma|\sigma^{\mathrm{co}})$ 的 Morita 理论构造，我们看到这一限制的像恰好等于 $\mathcal D_0^{-}$。

\subsubsection*{7.5.5.\quad $\mathrm{PGL}_2(\Q_p)$ 情形中 $\mathfrak A$ 的本质像}
我们现在考虑 $\IndCoh(\mathcal U)$ 的半正交分解沿 $j:\mathcal U\to\XX$ 的前推（一如既往地忽略非一般的成分，从而忽略幂零奇异支撑的问题，见附录 A.8.13）。由 (A.8.14)，我们有 $\IndCoh(\XX)^{-}$ 的半正交分解 $\mathcal A_1,\mathcal A_2$，其中 $\mathcal A_2=j_*\mathcal E_0^{-}$，而 $\mathcal A_1=(j^*)^{-1}(\mathcal D_0^{-})$。

我们断言 $\mathcal A_1$ 是 $\mathfrak A$ 的本质像。要看出这一点，注意对 $\mathcal A_1$ 中任何 $x$，我们有态射
\[
x\to j_*j^*x
\]
（来自 $j_*$ 与 $j^*$ 的伴随的单位），其余纤维（cofiber）支撑在 $\mathcal U$ 的补集上。该余纤维因此落在 $\mathfrak A$ 的本质像中，因为这个本质像在超奇异点处包含（奇偶性的）摩天层（它们是超奇异不可约表示的像，见定理 7.3.5）。因此只需证明：\ysrc{304}$\mathfrak A$ 的本质像在 $\mathcal U$ 上的限制正是 $\mathcal D_0^{-}$；而正如我们已经观察到的，这由 §7.3.20 中对该限制的显式描述可得。

\subsubsection*{7.5.6.\quad $D^\times$ 的表示}
为简单起见，我们继续假定 $p\ge5$。

我们先描述猜想性函子 $\mathfrak A_{D^\times}$ 的若干特征：它把中心特征平凡的光滑 $D^\times$-表示取到偶范畴 $\IndCoh(\XX)^{+}$。我们预期这样的函子存在，并且享有与我们关于 $\mathrm{PGL}_2(\Q_p)$ 的函子 $\mathfrak A$ 类似的性质。特别地，我们预期它是连续且全忠实的，并预期它满足与 Taylor--Wiles 拼接的相容性。

利用后一预期，我们可以预测 $\mathfrak A_{D^\times}$ 在一组生成元上的取值。为此我们简要回忆 $D^\times$ 的一些基本表示论。为避免讨论有理性的问题，我们以 $\ol{\F}_p$-系数工作。把 $D$ 的素元记作 $\Pi$，把极大序记作 $\mathcal O_D$，则 $\mathcal O_D/\Pi$ 是 $\F_p$ 的二次扩张，我们把它与 $\F_{p^2}$ 等同。于是 $\mathcal O_D^\times$ 的不可约 $\ol{\F}_p$-表示由 $\F_{p^2}^\times$ 的特征给出。固定嵌入 $\F_{p^2}\into\ol{\F}_p$ 便决定了特征 $\ol{\eta}:\mathcal O_D^\times\to\ol{\F}_p^\times$，而其他每个这样的特征都是某个 $i\in\Z/(p^2-1)\Z$ 对应的 $\ol{\eta}^{i}$ 的形式。这些表示的中心特征平凡，当且仅当 $(p-1)$ 整除 $i$；下文我们总假定如此。此时我们可以把它们延拓为 $\Pi^{2\Z}\mathcal O_D^\times$ 的、中心特征平凡的表示，我们仍记作 $\ol{\eta}^{i}$。

若 $(p+1)\nmid i$，则 $\sigma_i:=\mathrm{Ind}_{\Pi^{2\Z}\mathcal O_D^\times}^{D^\times}\ol{\eta}^{i}$ 是 $D^\times$ 的不可约二维表示（中心特征平凡）；我们有 $\sigma_i\cong\sigma_{-i}$，且当 $i\ne\pm j$ 时 $\sigma_i\not\cong\sigma_j$。若 $i=0$ 或 $i=(p^2-1)/2$，则 $\mathrm{Ind}_{\Pi^{2\Z}\mathcal O_D^\times}^{D^\times}\ol{\eta}^{i}$ 是两个特征 $\sigma_i^{\pm}$ 的直和，满足 $\sigma_i^{\pm}(\Pi)=\pm1$。这些表示给出了「中心特征平凡的光滑 $D^\times$-表示」这一范畴的一组生成元。

我们预期 $\mathfrak A_{D^\times}$ 在这些表示上的取值具有特别简单的形式。T.G. 与 David Geraghty 的模 $p$ Jacquet--Langlands 对应 [GG15] 把 $\ol{\eta}^{i}$ 对应到一对 Serre 权 $\sigma,\sigma^{\mathrm{co}}$ 的伴随对（定义 7.2.14 意义下）。事实上，由注 7.2.15，对 $(\sigma,\sigma^{\mathrm{co}})$ 正是某个尖点驯顺型的 Jordan--Hölder 成分之集；而这个尖点驯顺型对应于（经局部类域论对应于 $\ol{\eta}^{i}\oplus\ol{\eta}^{-i}$ 的 $I_{\Q_p}$ 之表示到特征零的驯顺提升）。

记 $\XX(\sigma_i):=\XX(\sigma)\cup\XX(\sigma^{\mathrm{co}})$。（注意这与定义 7.2.16 中的 $\XX(\sigma|\sigma^{\mathrm{co}})$ 并不总一致，因为我们略去了形如 $\XX(\sigma_{a,p-1})^{\pm}$ 的成分。）于是我们预期（$j=0$ 或 $1$）
\begin{equation}\tag{7.5.7}\label{eq:a4-7-5-7}
\mathfrak A_{D^\times}\bigl(\sigma^{\pm}_{(p^2-1)j/2}\bigr)
=\OO_{\XX(\sigma_{(p-1)j/2,p-1})^{\pm}},
\end{equation}
而当 $(p+1)\nmid i$ 时我们预期
\begin{equation}\tag{7.5.8}\label{eq:a4-7-5-8}
\mathfrak A_{D^\times}(\sigma_i)=\OO_{\XX(\sigma_i)}.
\end{equation}

我们对这一预期的论证如下。由模 $p$ 四元数模形式的重数一（multiplicity one）结果（即相应的拼接模的无挠性结果，见 §7.5.13），以及 $\mathfrak A_{D^\times}$ 与 Taylor--Wiles 拼接的预期相容性，我们预期所有这些层在其支撑（即所指的闭子叠）上都是可逆层。（据我们所知，\ysrc{305}这些重数一结果尚不见于文献，故我们在下文略述其证明。）由于它们也落在 $\IndCoh(\XX)^{+}$ 中，一个自然的猜测是：它们由各自支撑的结构层给出。

\begin{rem}{7.5.9}\label{thm:a4-7-5-9}
看起来合理的是：函子 $\mathfrak A_{D^\times}$ 与猜想 9.3.2 在 Eisenstein 情形下的相容性（对 $\Q$ 上在 $p$ 处分歧的四元数代数）会蕴含这些可逆层确实就是其支撑的结构层，从而为我们的预期提供进一步证据。
\end{rem}

\begin{rem}{7.5.10}\label{thm:a4-7-5-10}
若放弃「表示具有平凡中心特征」这一假设，则上述陈述以显然的方式推广：只要规定 Serre 权 $\sigma_{a,p-2}$ 是它自己的伴随（故此时「伴随对」是单点集）。
\end{rem}

现在我们可以重复 §7.5.5 的分析。我们再次有 $\IndCoh(\XX)^{+}$ 的半正交分解 $\mathcal A_1,\mathcal A_2$，其中 $\mathcal A_2=j_*\mathcal E_0^{+}$，而 $\mathcal A_1=(j^*)^{-1}(\mathcal D_0^{+})$。利用预期 (7.5.7) 与 (7.5.8)，我们看到 $\mathfrak A_{D^\times}$ 的本质像包含在 $\mathcal A_1$ 中，且 $\mathcal A_1$ 由该本质像与不可约点处的（偶）摩天层共同生成。

然而，与 $\mathrm{GL}_2(\Q_p)$ 的情形相反，$\mathfrak A_{D^\times}$ 的本质像并不包含这些摩天层！相反，我们预期这些摩天层由 $\Bun_{\mathrm{PGL}_2}$ 的第一个非开分层片——即由 $\mathrm{BC}(\OO(1))$ 的表示论——所解释。因此我们是在暗示 $\mathrm{BC}(\OO(1))$ 的行为与 $\mathrm{BC}(\OO(n))$（$n>1$）不同；对此我们指出一些理由：$\mathrm{BC}(\OO(n))$ 当 $n=1$ 时可由完满空间表示，而当 $n>1$ 时不能。\tn{底本此处作「these skyscraper accounts are accounted for by…」，其中 accounts 疑为 sheaves 之误；译文按上下文译作「摩天层」。}

\subsubsection*{7.5.11.\quad $\mathrm{GL}_2(\Q_{p^2})$ 的情形}
对 $\mathrm{GL}_2(\Q_{p^f})$ 的情形可以作类似分析。我们可以用与上述关于 $D^\times$ 的预测类似的方式，给出假想函子 $\mathfrak A$ 的本质像的生成元；并且我们可以描述相应叠 $\XX_{\Q_{p^f},2}$ 的不可约成分（对一般的 Serre 权，我们将在 §7.6 做这件事）。人们也可以通过类似于上述的分析，在可约轨迹上构造半正交分解。

$\mathrm{GL}_2(\Q_{p^2})$ 的情形与上述 $D^\times$ 的情形非常相似。我们再次发现（假想的）$\mathfrak A$ 的本质像包含在相应层范畴的偶部分中，并且包含在某个显式层范畴的左正交补中，该显式层范畴由从可约轨迹推出的层生成。我们再次发现不可约点处的摩天层包含在这个左正交补中，但似乎不在 $\mathfrak A$ 的本质像中。

我们再次预期这些摩天层由 Banach--Colmez 空间的表示论所解释，此情形下即 $\mathrm{BC}(\OO(2))$；同样，这个 Banach--Colmez 空间可由完满空间表示。（对一般的 $F$，$\mathrm{BC}(\OO(n))$ 可由完满空间表示当且仅当 $0<n\le[F:\Q_p]$，见 [FS24, Prop. II.2.5(iv)]。）

\begin{rem}{7.5.12}\label{thm:a4-7-5-12}
不可约点处的摩天层不落在 $\mathfrak A$ 的本质像中，这与 $\mathrm{GL}_2(\Q_{p^2})$ 的超奇异表示不能有限呈现有关；见下文 §7.7.59。我们不知道这种有限呈现性的失效是否有纯粹表示论上的解释，即用 $\mathrm{BC}(\OO(2))$ 来解释。\tn{底本此处引「§7.7.59」，与该节段落编号的体例不合，疑为笔误；译文照录。}
\end{rem}

\ysrc{306}
\subsubsection*{7.5.13.\quad 重数一}
本节最后，我们略述上面提到的关于 $D^\times$ 的重数一结果的证明。为完整起见，我们放弃「表示具有平凡中心特征」这一假设。设 $D_\Q$ 是 $\Q$ 上的四元数代数，恰在 $p,\infty$ 处分歧；并设 $\ol{r}:\Gal_\Q\to\mathrm{GL}_2(\ol{\F}_p)$ 使得 $\ol{r}|_{\Gal_{\Q_p}}$ 不可约。设 $\frakm$ 是某个适当的、与 $p$ 互素的 Hecke 代数的对应极大理想。（为简单起见，读者可以设想 $\ol{r}$ 在 $p$ 之外非分歧。）对 $\mathcal O_D^\times$ 的每个特征 $\ol{\eta}^{i}$，我们有相应的四元数模形式空间 $M(\ol{\eta}^{i})$（对群 $D_\Q^\times$），以及拼接模 $M_\infty(\ol{\eta}^{i})$——它是拼接模 $M_\infty(\eta^{i})$ 的模 $p$ 约化，其中 $\eta$ 是 $\ol{\eta}$ 到 $\F_{p^2}^\times$ 的 $\OO^\times$-值特征的提升。

这些拼接模已在 [GS11, GG15, Roz15] 等文中研究；由模性提升定理可以证明 $M_\infty(\eta^{i})$ 是某个环 $R_\infty(\eta^{i})$ 上的极大 Cohen--Macaulay 模。环 $R_\infty(\eta^{i})$ 在加框架形变环 $R_{\ol{\rho}}(\eta^{i})$ 上形式光滑，其中 $\ol{\rho}:=\ol{r}|_{\Gal_{\Q_p}}$；后者可描述如下。

首先，若 $(p+1)\mid i$，则扭之后可设 $i=0$，此时 $R_{\ol{\rho}}(\eta^{i})$ 是非晶体半稳定表示（Hodge--Tate 权为 $0,1$）的加框架形变环。其次，若 $(p+1)\nmid i$，则 $R_{\ol{\rho}}(\eta^{i})$ 是潜在晶体表示（Hodge--Tate 权为 $0,1$、惯性型为 $\eta^{i}\oplus\eta^{-i}$）的加框架形变环。

有了这一描述，容易验证拼接模的支撑与 (7.5.8)、(7.5.7) 的规定相符。（这是四元数模形式的 Serre 猜想的权部分，最初由 T.G. 与 Savitt 在 [GS11] 中证明；区分 $\sigma_i^{\pm}$ 的加细由 Rozensztajn 在 [Roz15, Lem. 6.1.1] 中证明，方法是考虑「扩张型」（extended types），即跟踪某个 Frobenius 的作用。）

现在，若 $R_{\ol{\rho}}(\eta^{i})$ 是正则的，则 $M_\infty(\eta^{i})$ 自动无挠（这是由 Taylor--Wiles 拼接推出重数一结果的通常论证，独立于 Diamond 与 Fujiwara）。特别地，只要 $R_{\ol{\rho}}(\eta^{i})$ 或者形式光滑，或者在 $\OO[[x,y]]/(xy-p)$ 上形式光滑，此结论就成立，其中 $\OO$ 是 $\Q_p$ 的某个适当非分歧扩大的整数环。形变环 $R_{\ol{\rho}}(\eta^{i})$ 已在许多情形下算出。在 $(p+1)\mid i$ 的情形，它们形式光滑（且可非常显式地描述，因为所有提升都是普通的）；而若 $(p+1)\nmid i$，则由 [Sav05, Thm. 1.4] 与 [EGS15, Thm. 7.2.1] 可知，除 $\ol{\rho}$ 是某个非分歧表示的扭、或 $\ol{\rho}|_{I_p}$ 是 $1\oplus\ol{\varepsilon}$ 的扭这两种情形外，它们都是正则的；下文我们把这些情形称为例外情形。

可以「纯思考」地证明（通过对比 Hodge--Tate 权为 $0,p-1$ 的晶体表示的拼接模）：若 $\ol{\rho}$ 是不同的非分歧特征之直和，则 $R_{\ol{\rho}}(\eta^{i})$ 并不正则。事实上，由 [LHMM24] 的结果，它在 $\OO[[x,y]]/(xy-p^2)$ 上形式光滑（至少当 $p$ 充分大时）。此情形下，似乎不太可能不用某些几何输入就证明 $M_\infty(\eta^{i})$ 的无挠性。其余情形下的形变环也在 [LHMM24] 中算出（对充分大的 $p$）；细节见 [LHMM24, §5.5]。

下述关于这些例外情形中无挠性的论证取自 George Boxer、Frank Calegari 与 T.G. 的未发表工作。（它由 Chengxiang Bao、Andrea Dotto 与 Yulia Kotelnikova 独立发现。）由 Nakayama 引理，只需证明相应的四元数模形式特征空间\ysrc{307}$M(\ol{\eta}^{i})[\frakm]$ 在 $\ol{\F}_p$ 上是一维。

\begin{rem}{7.6.1}\label{thm:a4-7-6-1}
熟悉 $\mathrm{GL}_2(\Q_{p^2})$ 的 $p$-进局部 Langlands 纲领中种种困难与神秘之处、并急于看到我们能对它们提供何种线索的读者，可以直接跳到 §7.7.59 与注 7.8.6。
\end{rem}

\subsection{$\mathrm{GL}_2(\Q_{p^f})$ 的 Banach 情形中的预备}
继 $\mathrm{GL}_2(\Q_p)$ 的情形之后，下一个自然要考虑的情形是 $G:=\mathrm{GL}_2(\Q_{p^f})$（$f>1$）；这里一如既往地 $\Q_{p^f}$ 表示 $\Q_p$ 的非分歧扩张，剩余域为 $\F_q:=\F_{p^f}$，并记 $K:=\mathrm{GL}_2(\Z_{p^f})$、$Z=Z(G)$。例如，此情形下的模 $p$ 局部 Langlands 对应已被相当深入地研究，模 $p$ 局部--整体相容性问题亦然。（我们在行文中会顺带回忆这些研究的部分历史；不熟悉的读者不妨参阅 [Bre10a] 以及 [BHH$^{+}$21] 的引言。）另一方面，模 $p$（更一般地，模 $p^n$）光滑 $G$-表示的范畴之结构仍然相当神秘，超奇异不可约容许表示的分类仍未知，而且至今还没有 $\mathrm{GL}_2(\Q_p)$ 情形中 [Col10b] 与 [Pas13] 所例示的那种「传统」模式的模 $p$ 或 $p$-进 Langlands 对应的精确表述。

我们相信，范畴的视角能为这一情形提供不少洞见，既有助于阐明已知结果，也有助于说明为什么 $p$-进 Langlands 对应的传统表述方式显得如此难以捉摸。在接下来的讨论中，我们通过描述我们\footnote{这些想法中的许多是在 M.E. 与 T.G. 同 Ana Caraiani、Michael Harris、Bao Le Hung、Brandon Levin 与 David Savitt 的合作中发展出来的，属于 NSF 资助 1952556「Collaborative Research: Geometric Structures in the p-adic Langlands Program」的支持范围；我们感谢他们允许我们在此描述这些结果。}关于 $\mathrm{GL}_2(\Q_{p^f})$ 的范畴 Langlands 纲领的若干预备结果与预期，来展开这些要点。为便于记号并明确想法，我们主要聚焦 $f=2$ 的情形，但本节开头先作一些对任何 $f$ 都成立的考量。

\begin{rem}{7.6.2}\label{thm:a4-7-6-2}
$\mathrm{GL}_2(\Q_p)$ 情形与 $\mathrm{GL}_2(\Q_{p^2})$ 情形的相似与差异，可由下例说明。两种情形下，对一般的 Serre 权 $\sigma$，函子 $\mathfrak A$ 都把 $c\text{-}\mathrm{Ind}_{KZ}^G\sigma$ 取到支撑在 $\XX_{\mathrm{red}}$ 的相应成分 $\XX(\sigma)$ 上的显式线丛。在 $\mathrm{GL}_2(\Q_p)$ 的情形，这一对应「不太导出」；特别地，简单计算表明：\ysrc{308}超奇异不可约容许表示 $(c\text{-}\mathrm{Ind}_{KZ}^G\sigma)/T$ 被送到支撑在相应不可约伽罗瓦表示处的（扭）摩天层。

对 $\mathrm{GL}_2(\Q_{p^2})$，同样的计算表明 $(c\text{-}\mathrm{Ind}_{KZ}^G\sigma)/T$ 被送到支撑在 $\XX(\sigma)$ 上对应于 $T=0$ 的闭轨迹处的层。然而这一支撑现在是一维的，而且表示 $(c\text{-}\mathrm{Ind}_{KZ}^G\sigma)/T$ 具有无限长度。事实上，我们知道（见 (7.7.53)）它把完整的紧诱导表示 $c\text{-}\mathrm{Ind}_{KZ}^G\sigma'$（对某些权 $\sigma'\ne\sigma$）作为子表示包含在内；由此容易推出：$\mathrm{GL}_2(\Q_{p^2})$ 的超奇异不可约容许表示不可能被取到真正的层（而非层的复形），这与 $\mathrm{GL}_2(\Q_p)$ 的情形（或 $\ell\ne p$ 的情形）形成鲜明对比。
\end{rem}

\subsubsection*{7.6.3.\quad $\XX_{\mathrm{red}}$ 的一般结构}
一如既往，我们固定系数 $\OO$ 为 $\Q_p$ 的某个有限扩张 $L$ 的整数环，并进一步假定存在 $\Q_{p^f}$ 到 $L$ 的嵌入。（换言之，我们可以把 $L$ 看作 $\Q_{p^f}$ 的扩张。）记 $k$ 为 $\OO$ 的剩余域。我们还固定一个特征 $\zeta:\Q_{p^f}^\times\to\OO^\times$，并记 $\XX$ 为分类「定义在 $p$-进完备 $\OO$-代数上、$\Gal_{\Q_{p^f}}$ 的秩 $2$ 投射 étale $(\varphi,\Gamma)$-模」的 Noetherian 形式代数叠，其行列式固定等于 $\zeta\varepsilon^{-1}$。

\subsubsection*{7.6.4.\quad 成分 $\XX(\sigma)$}
正如 $\mathrm{GL}_2(\Q_p)$ 的情形，底层的约化代数子叠 $\XX_{\mathrm{red}}$ 在 $\Spec k$ 上有限呈现，并且是有限多个不可约成分 $\XX(\sigma)$ 的并，这些成分由中心特征为 $\zeta|_{\OO^\times}$ 的 Serre 权 $\sigma$ 标号；唯一的例外是：当 $\sigma$ 是 Steinberg 权的扭时，出现一对成分 $X(\sigma)^{\pm}$\tn{底本此处写作直立体 $X(\sigma)^{\pm}$，与他处花体的 $\XX(\sigma)$ 不一致；译文照录。}。并且一如既往，对一般的 $\sigma$（形如
\[
\sigma=(\Sym^{a_0}\otimes\det^{b_0})\otimes\cdots\otimes(\Sym^{a_i}\otimes\det^{b_i})^{\mathrm{Fr}^{i}}\otimes\cdots\otimes(\Sym^{a_{f-1}}\otimes\det^{b_{f-1}})^{\mathrm{Fr}^{f-1}}
\]
且所有 $a_i$ 满足 $0\le a_i\le p-3$ 者），$\XX(\sigma)$ 可描述为系数在 $k$-代数中的 Fontaine--Laffaille 模的模叠。

一般地，对 $\F_p$ 上的约化群 $G$，相应的 Fontaine--Laffaille 模空间可描述为商 $(U\backslash G)/T$，其中 $U$ 是 $G$ 中某个固定 Borel 的幂幺根，经左平移作用，而 $T$ 是同一 Borel 的极大环面，经 Frobenius 扭共轭作用。在我们的情形，群 $G$ 是 $\mathrm{SL}_2$ 从 $\F_{p^f}$ 到 $\F_p$ 的纯量限制。由于假定 $k$ 包含 $\F_{p^f}$，我们有
\[
G=G_{/k}=\mathrm{SL}_2^{f},
\]
\[
U=U_{/k}=\left(\begin{pmatrix}1&\mathbf G_a\\0&1\end{pmatrix}\right)^{f},
\]
以及
\[
T=T_{/k}=\mathbf G_m^{f},
\]
Frobenius 循环置换各因子。

现在 $\begin{pmatrix}1&\mathbf G_a\\0&1\end{pmatrix}\backslash\mathrm{SL}_2=\mathbf A^{2}\setminus\{0\}$（经 $\begin{pmatrix}x&*\\*&y\end{pmatrix}\mapsto(x,y)$）。于是我们有
\[
U\backslash G:=\prod_{i=0}^{f-1}(\mathbf A^{2}\setminus\{0\})_{(x_i,y_i)}
\]
\ysrc{309}（下标表示我们对坐标的记法）。环面 $T$ 随后经
\[
(t_0,\dots,t_{f-1})\cdot(x_i,y_i)=(t_{i-1}t_ix_i,t_{i-1}t_i^{-1}y_i)
\]
作用。于是（对一般的 $\sigma$）我们发现 $\XX(\sigma)$ 等于商叠
\[
\XX(\sigma)=\Bigl[\Bigl(\prod_{i=0}^{f-1}(\mathbf A^{2}\setminus\{0\})_{(x_i,y_i)}\Bigr)/\mathbf G_m^{f}\Bigr],
\]
其中环面作用由前述公式给出。

$\XX(\sigma)$ 上的坐标 $(x_i,y_i)$ 显然并不完全典范，但也并非完全没有意义；特别地，我们很快会看到，它们的零点轨迹描述 $\XX(\sigma)$ 与其他成分相交的轨迹。乘积
\[
T:=y_0\cdots y_{f-1}
\]
起特殊作用。显然 $T$ 是 $\mathbf G_m^{f}$-不变的，因而在 $\XX(\sigma)$ 上定义一个正则函数。事实上它生成 $\XX(\sigma)$ 上的函数环，我们把它记录为下述引理。

\begin{lem}{7.6.5}\label{thm:a4-7-6-5}
\begin{enumerate}[label=(\arabic*),leftmargin=2.4em,itemsep=0.15em]
\item $\Gamma(\XX(\sigma),\OO_{\XX(\sigma)})=k[T]$。
\item 可区分的开子集 $D(T)$（即 $T$ 的非零轨迹）在 $\XX(\sigma)$ 中恰是 niveau 1 普通轨迹（niveau 1 ordinary locus）。
\end{enumerate}
\end{lem}

具体地，我们看到
\[
D(T)=[\mathbf A^{f}_{(x_0,\dots,x_{f-1})}/(\mathbf G_m)_t]\times(\mathbf G_m)_T,
\]
其中 $t$ 经
\[
t\cdot(x_0,\dots,x_{f-1})=(t^2x_0,\dots,t^2x_{f-1})
\]
作用；因此 $(\mathbf G_m)_t$ 可与 $\mathbf G_m$ 在 $\mathbf G_m^{f}$ 中的对角拷贝等同。存在特征 $\chi:\Gal_{\Q_{p^f}}\to k^\times$（依赖于 $\sigma$），使得 $D(T)$ 的点 $(T,x_0,\dots,x_{f-1})$ 对应于 $\Gal_{\Q_{p^f}}$ 的一个形如
\[
\begin{pmatrix}\mathrm{ur}_T\chi&*\\0&\mathrm{ur}_{T^{-1}}\zeta\varepsilon^{-1}\chi^{-1}\end{pmatrix}
\]
的二维表示 $\ol{\rho}$，其中 $*$ 是 $\Ext^{1}_{\Gal_{\Q_{p^f}}}(\mathrm{ur}_{T^{-1}}\zeta\varepsilon^{-1}\chi^{-1},\mathrm{ur}_T\chi)$ 的一个元素，而该 $\Ext$ 群是 $f$ 维向量空间，坐标为 $x_0,\dots,x_{f-1}$。

零点轨迹 $Z(T):=\XX(\sigma)\setminus D(T)$ 的几何更为复杂；我们在 §7.7 中详细描述 $f=2$ 的情形。

现设 $\sigma$ 非常一般，即每个 $a_i$ 满足 $0<a_i<p-3$。则有 $2f$ 个 Serre 权 $\sigma'$ 使得 $\Ext^{1}_{\mathrm{GL}_2(\F_q)}(\sigma,\sigma')\ne0$，且每个这样的 $\sigma'$ 都是一般的。此时可发现，$\XX(\sigma)$ 与 $\XX(\sigma')$ 恰在 $\XX(\sigma)$ 的一个余维 $1$ 的闭子叠中相交，该闭子叠由 $x_i=0$ 或 $y_i=0$（对某个 $i$）切出。

\ysrc{310}
\subsubsection*{7.6.6.\quad Picard 群}
继续假定 $\sigma$ 一般。由 $\XX(\sigma)$ 作为商叠的描述，容易看出 $\mathrm{Pic}(\XX(\sigma))\cong\Z^{f}$（$\mathbf G_m^{f}$ 的特征标格）。这一同构并不显然典范；但下面我们说明如何使它典范到 $\Z^{f}$ 的坐标置换，附加假设是 $\sigma$ 非常一般。

首先考虑 $\XX(\sigma)$ 的开子叠 $D(T)$。如前所述，它容许描述 $D(T)=(\mathbf A^{f}/\mathbf G_m)\times\mathbf G_m$，由此推出 $\mathrm{Pic}(D(T))=\Z$（$\mathbf G_m$ 的特征标格）。此外，这一等同是典范的，因为出现的 $\mathbf G_m$ 自然地有定向（这对应于如下事实：$D(T)$ 参数化沿固定方向由 $\mathrm{ur}_T\chi$ 扩张 $\mathrm{ur}_{T^{-1}}\zeta\varepsilon^{-1}\chi^{-1}$ 所得的扩张）。

对每个 $i=0,\dots,f-1$，记 $\sigma_i$ 为使 $\XX(\sigma)\cap\XX(\sigma_i)=\XX(\sigma)^{y_i=0}$ 的权；由于 $\sigma$ 非常一般，每个 $\sigma_i$ 也一般。记 $T_i$ 为 $\XX(\sigma_i)$ 上函数 $T$ 的类似物；则
\[
(T_i)|_{\XX(\sigma_i)^{y_i=0}}=\frac{x_{i-1}}{x_i}\prod_{j\ne i-1,i}y_j.
\]
因此
\[
\XX(\sigma)^{y_i=0}\cap D(T_i)=\XX(\sigma)^{y_i=0,\,(x_{i-1}\prod_{j\ne i-1,i}y_j)\ne0}.
\]
现在
\[
D\Bigl(x_{i-1}x_i\prod_{j\ne i-1,i}y_j\Bigr)
=[\mathbf A^{f}_{(x_0,\dots,x_{i-2},y_{i-1},y_i,x_{i+1},\dots,x_{f-1})}/(\mathbf G_m)_t],
\]
其中 $t$ 经
\begin{align*}
&t\cdot(x_0,\dots,x_{i-2},y_{i-1},y_i,x_{i+1},\dots,x_{f-1})\\
&\qquad=(t^2x_0,\dots,t^2x_{i-2},t^2y_{i-1},t^{-2}y_i,t^2x_{i+1},\dots,t^2x_{f-1})
\end{align*}
作用；环面 $(\mathbf G_m)_t$ 可看作 $\mathbf G_m$ 经
\[
t\mapsto(t,\dots,t,t^{-1}\text{ 位于第 }i-1\text{ 个槽位},t,\dots,t)
\]
嵌入 $\mathbf G_m^{f}$ 所得的拷贝。

闭子叠 $\XX(\sigma)^{y_i=0}\cap D(T_i)$ 随后是 $[\mathbf A^{f-1}/\mathbf G_m]$ 被 $y_i=0$ 切出的拷贝。当然，嵌入
\[
\XX(\sigma)^{y_i=0}\cap D(T_i)\into D(T_i)
\]
把 $[\mathbf A^{f-1}/\mathbf G_m]$ 嵌入 $[\mathbf A^{f}/\mathbf G_m]$ 的另一个拷贝。

我们已经观察到 $\mathrm{Pic}(D(T_i))=\mathrm{Pic}([\mathbf A^{f}/\mathbf G_m])=\Z$，且限制映射 $\mathrm{Pic}(D(T_i))\to\mathrm{Pic}([\mathbf A^{f-1}/\mathbf G_m])=\Z$ 是同构。因此 $\mathrm{Pic}(\XX(\sigma)^{y_i=0}\cap D(T_i))=\Z$（典范地）。现在考虑这些限制态射的乘积
\begin{equation}\tag{7.6.7}\label{eq:a4-7-6-7}
\begin{aligned}
\Z^{f}\cong\mathrm{Pic}(\XX(\sigma))&\to\mathrm{Pic}(D(T))\times\prod_{i=0}^{f-1}\mathrm{Pic}(\XX(\sigma)^{y_i=0}\cap D(T_i))\\
&=\Z\times\Z^{f}=\Z^{f+1}.
\end{aligned}
\end{equation}
回顾 $D(T)$ 与各个 $\XX(\sigma)^{y_i=0}\cap D(T_i)$ 作为商 $[\mathbf A^{f}/(\mathbf G_m)_t]$ 与 $[\mathbf A^{f-1}/(\mathbf G_m)_t]$ 的描述，以及这些各个拷贝

% 本片止于印刷页 310 页末（§7.6.6）。上句在底本 PDF 324 页（印刷页 311）继续：
% "…of (G_m)_t into (G_m)^f, we find that this map is given by [公式]"，续片自 PDF 324 页起接「的嵌入……」续译。
